\documentclass[10pt,a4paper]{article}

% ── Packages ──────────────────────────────────────────────────────────────────
\usepackage[top=1.8cm,bottom=1.8cm,left=1.6cm,right=1.6cm]{geometry}
\usepackage{fontenc}
\usepackage[utf8]{inputenc}
\usepackage{microtype}
\usepackage{parskip}
\usepackage{titlesec}
\usepackage{titletoc}
\usepackage{booktabs}
\usepackage{tabularx}
\usepackage{array}
\usepackage{multirow}
\usepackage{xcolor}
\usepackage{tcolorbox}
\usepackage{enumitem}
\usepackage{amsmath,amssymb}
\usepackage{multicol}
\usepackage{fancyhdr}
\usepackage{hyperref}
\usepackage{mdframed}
\usepackage{graphicx}
\usepackage{colortbl}
\usepackage{makecell}

% ── Colors ────────────────────────────────────────────────────────────────────
\definecolor{NavyBlue}{RGB}{25,55,110}
\definecolor{DarkBlue}{RGB}{10,40,90}
\definecolor{LightBlue}{RGB}{220,232,247}
\definecolor{MedBlue}{RGB}{180,210,240}
\definecolor{GreenTeal}{RGB}{0,110,90}
\definecolor{LightGreen}{RGB}{220,245,235}
\definecolor{WarnRed}{RGB}{160,30,30}
\definecolor{LightRed}{RGB}{250,225,225}
\definecolor{AmberOrange}{RGB}{180,90,0}
\definecolor{LightAmber}{RGB}{255,240,210}
\definecolor{PurpleDeep}{RGB}{80,0,120}
\definecolor{LightPurple}{RGB}{240,225,255}
\definecolor{RowGray}{RGB}{245,245,248}
\definecolor{HeaderGray}{RGB}{60,60,75}

% ── tcolorbox styles ──────────────────────────────────────────────────────────
\tcbuselibrary{skins,breakable}
\tcbset{
  when/.style={
    colback=LightBlue, colframe=NavyBlue, coltitle=white,
    fonttitle=\bfseries\small, title={When to Use},
    breakable, left=4pt, right=4pt, top=3pt, bottom=3pt,
    arc=2pt, boxrule=0.8pt
  },
  tests/.style={
    colback=LightGreen, colframe=GreenTeal, coltitle=white,
    fonttitle=\bfseries\small, title={Diagnostic Tests \& Decision Rules},
    breakable, left=4pt, right=4pt, top=3pt, bottom=3pt,
    arc=2pt, boxrule=0.8pt
  },
  watchout/.style={
    colback=LightRed, colframe=WarnRed, coltitle=white,
    fonttitle=\bfseries\small, title={Watch Out For},
    breakable, left=4pt, right=4pt, top=3pt, bottom=3pt,
    arc=2pt, boxrule=0.8pt
  },
  checklist/.style={
    colback=LightAmber, colframe=AmberOrange, coltitle=white,
    fonttitle=\bfseries\small, title={Researcher Checklist},
    breakable, left=4pt, right=4pt, top=3pt, bottom=3pt,
    arc=2pt, boxrule=0.8pt
  },
  key/.style={
    colback=LightPurple, colframe=PurpleDeep, coltitle=white,
    fonttitle=\bfseries\small, title={Key Formula / Estimator},
    breakable, left=4pt, right=4pt, top=3pt, bottom=3pt,
    arc=2pt, boxrule=0.8pt
  }
}

% ── Section styles ────────────────────────────────────────────────────────────
\titleformat{\section}
  {\color{DarkBlue}\large\bfseries\sffamily}
  {\color{DarkBlue}\thesection.}{6pt}{}
  [\color{NavyBlue}\titlerule]

\titleformat{\subsection}
  {\color{HeaderGray}\normalsize\bfseries\sffamily}
  {\thesubsection.}{5pt}{}

\titlespacing{\section}{0pt}{10pt}{4pt}
\titlespacing{\subsection}{0pt}{6pt}{2pt}

% ── List settings ─────────────────────────────────────────────────────────────
\setlist[itemize]{leftmargin=1.2em, topsep=1pt, itemsep=0pt, parsep=0pt}
\setlist[enumerate]{leftmargin=1.4em, topsep=1pt, itemsep=0pt, parsep=0pt}
\newlist{checklist}{itemize}{1}
\setlist[checklist]{label=$\square$, leftmargin=1.4em,
  topsep=1pt, itemsep=1pt, parsep=0pt}

% ── Header / footer ───────────────────────────────────────────────────────────
\pagestyle{fancy}
\fancyhf{}
\fancyhead[L]{\small\color{NavyBlue}\bfseries Applied Econometrics — Model Selection \& Diagnostic Checklist}
\fancyhead[R]{\small\color{NavyBlue} University of Neuch\^{a}tel}
\fancyfoot[C]{\small\color{gray}\thepage}
\renewcommand{\headrulewidth}{0.6pt}
\renewcommand{\headrule}{\color{NavyBlue}\hrule width\headwidth height\headrulewidth}

% ── Hyperref ──────────────────────────────────────────────────────────────────
\hypersetup{
  colorlinks=true, linkcolor=NavyBlue, urlcolor=GreenTeal,
  pdftitle={Applied Econometrics Reference Guide},
  pdfauthor={University of Neuchatel}
}

% ── Helpers ───────────────────────────────────────────────────────────────────
\newcommand{\pval}[1]{\textit{p}-value #1 0.05}
\newcommand{\reject}{$\Rightarrow$ \textbf{reject H\textsubscript{0}}}
\newcommand{\fail}{$\Rightarrow$ \textbf{fail to reject H\textsubscript{0}}}
\newcommand{\pkg}[1]{\texttt{#1}}
\newcommand{\fn}[1]{\texttt{#1}}

% ══════════════════════════════════════════════════════════════════════════════
\begin{document}
% ══════════════════════════════════════════════════════════════════════════════

% ── Title page ────────────────────────────────────────────────────────────────
\begin{titlepage}
\pagecolor{NavyBlue}
\color{white}
\vspace*{3cm}
\begin{center}
{\fontsize{28}{34}\selectfont\bfseries\sffamily Applied Econometrics}\\[0.5em]
{\fontsize{18}{24}\selectfont\sffamily Model Selection \& Diagnostic Checklist}\\[1.5em]
\textcolor{MedBlue}{\rule{12cm}{1pt}}\\[1.5em]
{\large\sffamily A Practical Reference for Researchers}\\[0.5em]
{\normalsize\sffamily\color{MedBlue}
When to use each model $\cdot$ Key assumptions $\cdot$ Diagnostic tests\\
Decision rules $\cdot$ Common pitfalls $\cdot$ Step-by-step checklists}
\vfill
{\large\sffamily University of Neuch\^{a}tel}\\[0.3em]
{\normalsize\sffamily\color{MedBlue} Python for Econometrics Course}\\[0.3em]
{\normalsize\sffamily\color{MedBlue} \today}
\end{center}
\end{titlepage}
\nopagecolor

% ── TOC ───────────────────────────────────────────────────────────────────────
\tableofcontents
\thispagestyle{fancy}
\newpage

% ══════════════════════════════════════════════════════════════════════════════
% QUICK DECISION TREE
% ══════════════════════════════════════════════════════════════════════════════
\section{Quick Model-Selection Decision Tree}

\begin{tcolorbox}[colback=LightBlue,colframe=NavyBlue,arc=3pt,boxrule=0.8pt]
\small
\textbf{Step 1 — What is your data structure?}
\begin{itemize}
  \item Cross-sectional (one observation per unit) $\to$ Sections 2–7
  \item Time series (one unit over time) $\to$ Sections 8–10
  \item Panel (multiple units over time) $\to$ Section 11
  \item Binary / limited outcome $\to$ Section 13
\end{itemize}

\textbf{Step 2 — Is the outcome continuous and unbounded?}\\
Yes $\to$ OLS family. \quad No (binary/count/censored) $\to$ Section 13.

\textbf{Step 3 — Is your key regressor endogenous?}\\
Suspect endogeneity (OVB, measurement error, simultaneity)?  $\to$ Section 12 (IV/2SLS) or Section 14 (Treatment Effects).

\textbf{Step 4 — Run OLS, then check diagnostics in order:}\\
Multicollinearity (Sec.~3) $\to$ Heteroscedasticity (Sec.~4) $\to$
Specification (Sec.~5) $\to$ Serial correlation if time series (Sec.~9) $\to$
Stationarity if time series (Sec.~8, 10).
\end{tcolorbox}

\bigskip

\begin{center}
\footnotesize
\begin{tabular}{lllll}
\toprule
\rowcolor{NavyBlue}\color{white}\textbf{Setting} &
\color{white}\textbf{First choice} &
\color{white}\textbf{If endogeneity} &
\color{white}\textbf{If binary $Y$} &
\color{white}\textbf{If non-stationary} \\
\midrule
Cross-section & OLS & IV / 2SLS & Logit / Probit & --- \\
\rowcolor{RowGray}
Time series & OLS + HAC & IV & --- & ARIMA / ECM \\
Panel & FE / RE & FE + IV & --- & --- \\
\rowcolor{RowGray}
Program eval. & DiD / RDD & IV-LATE & --- & --- \\
\bottomrule
\end{tabular}
\end{center}

\newpage

% ══════════════════════════════════════════════════════════════════════════════
\section{OLS --- Simple \& Multiple Regression}
% ══════════════════════════════════════════════════════════════════════════════

\begin{multicols}{2}

\begin{tcolorbox}[when]
\begin{itemize}
  \item Continuous, unbounded dependent variable
  \item Exploring linear relationships between variables
  \item Baseline before testing for violations
  \item Gauss--Markov conditions plausibly satisfied
\end{itemize}
\end{tcolorbox}

\begin{tcolorbox}[key]
\[
\hat{\boldsymbol{\beta}} = (\mathbf{X}'\mathbf{X})^{-1}\mathbf{X}'\mathbf{y}
\]
\[
\text{Var}(\hat{\boldsymbol{\beta}}) = \sigma^2(\mathbf{X}'\mathbf{X})^{-1}
\]
BLUE when MLR.1--MLR.5 hold (Gauss--Markov Theorem).
\end{tcolorbox}

\columnbreak

\begin{tcolorbox}[watchout]
\begin{itemize}
  \item Omitted variable bias: $\text{Bias} = \beta_2 \cdot \delta_1$
  \item Functional form misspecification (try logs, quadratics)
  \item Influential outliers (check Cook's distance)
  \item Extrapolation beyond data support
  \item Interpreting $R^2$ as causal fit
\end{itemize}
\end{tcolorbox}

\end{multicols}

\begin{tcolorbox}[tests]
\begin{tabularx}{\linewidth}{lXl}
\toprule
\textbf{Test} & \textbf{Null hypothesis} & \textbf{Decision} \\
\midrule
$t$-test & $H_0: \beta_j = 0$ & $|t| > t_{\alpha/2,n-k-1}$ \reject \\
\rowcolor{RowGray}
$F$-test (joint) & $H_0:$ all slope $\beta = 0$ & $F > F_{\alpha,k,n-k-1}$ \reject \\
RESET test & $H_0:$ functional form OK & \pval{<} $\Rightarrow$ misspecification \\
\rowcolor{RowGray}
Partial $F$-test & $H_0:$ group of $\beta$s $= 0$ & $F > F_{\alpha,q,n-k-1}$ \reject \\
\bottomrule
\end{tabularx}
\end{tcolorbox}

\begin{tcolorbox}[checklist]
\begin{multicols}{2}
\begin{checklist}
  \item State the population model explicitly
  \item Check variable distributions (histograms)
  \item Consider log transformation for skewed vars
  \item Run OLS, inspect coefficients and signs
  \item Check $R^2$ and adjusted $R^2$
  \item Run RESET test for functional form
  \item Inspect residual plot vs.\ fitted values
  \item Run joint $F$-test for overall significance
  \item Check for influential observations
  \item Interpret coefficients as \textit{ceteris paribus}
\end{checklist}
\end{multicols}
\end{tcolorbox}

% ══════════════════════════════════════════════════════════════════════════════
\section{Multicollinearity}
% ══════════════════════════════════════════════════════════════════════════════

\begin{multicols}{2}

\begin{tcolorbox}[when]
\begin{itemize}
  \item After OLS, if SEs are large / signs unexpected
  \item Multiple regressors that may be highly correlated
  \item Model includes polynomial or interaction terms
  \item Structural economic variables (income, wealth, consumption)
\end{itemize}
\end{tcolorbox}

\begin{tcolorbox}[key]
\[
\text{VIF}_j = \frac{1}{1 - R^2_j}
\]
where $R^2_j$ = $R^2$ from regressing $x_j$ on all other regressors.\\
Condition number: $\kappa = \sqrt{\lambda_{\max}/\lambda_{\min}}$
\end{tcolorbox}

\columnbreak

\begin{tcolorbox}[watchout]
\begin{itemize}
  \item Inflated SEs $\to$ false non-significance
  \item Coefficients sensitive to adding/dropping vars
  \item Signs may flip relative to bivariate correlations
  \item High $R^2$ with few significant $t$-tests
  \item \textbf{Not a violation} of Gauss--Markov; OLS still BLUE
\end{itemize}
\end{tcolorbox}

\end{multicols}

\begin{tcolorbox}[tests]
\begin{tabularx}{\linewidth}{lXl}
\toprule
\textbf{Indicator} & \textbf{What it measures} & \textbf{Threshold} \\
\midrule
VIF$_j$ & Variance inflation for $\hat\beta_j$ & $>10$ severe, $>5$ moderate \\
\rowcolor{RowGray}
Condition number $\kappa$ & Overall collinearity & $>30$ severe \\
Pairwise $r$ & Simple correlation & $|r| > 0.8$ flag \\
\rowcolor{RowGray}
Eigenvalue ratio & Near-zero eigenvalues & $\lambda < 0.01$ flag \\
\bottomrule
\end{tabularx}
\end{tcolorbox}

\begin{tcolorbox}[checklist]
\begin{multicols}{2}
\begin{checklist}
  \item Compute pairwise correlations of regressors
  \item Compute VIF for each regressor
  \item Check condition number of $\mathbf{X}'\mathbf{X}$
  \item If VIF $>$ 10: consider dropping one variable
  \item Consider combining correlated vars (index/PCA)
  \item Try Ridge if prediction is the goal
  \item Report VIFs in results table
  \item Do not drop variables with theoretical justification
\end{checklist}
\end{multicols}
\end{tcolorbox}

\newpage

% ══════════════════════════════════════════════════════════════════════════════
\section{Heteroscedasticity}
% ══════════════════════════════════════════════════════════════════════════════

\begin{multicols}{2}

\begin{tcolorbox}[when]
\begin{itemize}
  \item Cross-sectional data with varying unit sizes (firms, countries)
  \item Income/expenditure data (variance rises with level)
  \item After OLS, residuals show fan/funnel pattern
  \item Any time you want valid inference from OLS
\end{itemize}
\end{tcolorbox}

\begin{tcolorbox}[key]
Heteroscedasticity-robust (HC1) SEs:
\[
\widehat{\text{Var}}_{\text{HC1}} = \frac{n}{n-k-1}(\mathbf{X}'\mathbf{X})^{-1}\!\left(\sum_i \hat{u}_i^2 \mathbf{x}_i\mathbf{x}_i'\right)\!(\mathbf{X}'\mathbf{X})^{-1}
\]
WLS: weight by $w_i = 1/\hat{\sigma}_i^2$.
\end{tcolorbox}

\columnbreak

\begin{tcolorbox}[watchout]
\begin{itemize}
  \item OLS coefficients still unbiased, SEs wrong
  \item $t$ and $F$ statistics invalid without correction
  \item White test has low power in small samples
  \item WLS requires correct variance specification
  \item Always use robust SEs by default in cross-section
\end{itemize}
\end{tcolorbox}

\end{multicols}

\begin{tcolorbox}[tests]
\begin{tabularx}{\linewidth}{lXll}
\toprule
\textbf{Test} & \textbf{$H_0$} & \textbf{Statistic} & \textbf{Decision} \\
\midrule
Breusch--Pagan & Homoscedastic & $nR^2 \sim \chi^2_k$ & $p < 0.05$ \reject \\
\rowcolor{RowGray}
White & Homoscedastic & $nR^2 \sim \chi^2_p$ & $p < 0.05$ \reject \\
Goldfeld--Quandt & Homoscedastic & $F$-ratio of variances & $p < 0.05$ \reject \\
\bottomrule
\end{tabularx}
\end{tcolorbox}

\textbf{Remedy hierarchy:} (1) Use HC1/HC3 robust SEs $\to$ (2) WLS if variance structure known $\to$ (3) Log-transform $y$ if level variance rises with level.

\begin{tcolorbox}[checklist]
\begin{multicols}{2}
\begin{checklist}
  \item Plot residuals vs.\ fitted values
  \item Plot $|\hat{u}|$ or $\hat{u}^2$ vs.\ each regressor
  \item Run Breusch--Pagan test
  \item Run White test (more general)
  \item If tests reject: use \texttt{cov\_type='HC1'}
  \item If variance structure known: use WLS
  \item Compare OLS vs.\ robust SEs in table
  \item Check if log-transform eliminates problem
\end{checklist}
\end{multicols}
\end{tcolorbox}

% ══════════════════════════════════════════════════════════════════════════════
\section{Model Specification \& Data Issues}
% ══════════════════════════════════════════════════════════════════════════════

\begin{multicols}{2}

\begin{tcolorbox}[when]
\begin{itemize}
  \item Always: before trusting OLS results
  \item When theory suggests non-linear relationships
  \item When proxy variables used (e.g., IQ for ability)
  \item When measurement error in key regressor suspected
\end{itemize}
\end{tcolorbox}

\begin{tcolorbox}[key]
\textbf{Attenuation bias (measurement error):}
\[
\text{plim}(\hat\beta_1) = \beta_1 \cdot \frac{\sigma^2_{x^*}}{\sigma^2_{x^*} + \sigma^2_v}
\]
\textbf{RESET statistic:} Add $\hat{y}^2, \hat{y}^3$ to regression; test their joint significance.
\end{tcolorbox}

\columnbreak

\begin{tcolorbox}[watchout]
\begin{itemize}
  \item Measurement error biases $\hat\beta$ toward zero
  \item Non-linearity looks like heteroscedasticity
  \item Proxy variables introduce attenuation
  \item Outliers can drive OLS results in small samples
  \item Missing variable vs.\ irrelevant variable trade-off
\end{itemize}
\end{tcolorbox}

\end{multicols}

\begin{tcolorbox}[checklist]
\begin{multicols}{2}
\begin{checklist}
  \item Run RESET test ($p < 0.05 \to$ misspecification)
  \item Try log, squared, interaction terms if RESET fails
  \item Check for outliers; run model with and without
  \item Document proxy variable limitations
  \item If meas.\ error: consider IV for the mismeasured var
  \item Davidson--MacKinnon $J$-test for non-nested models
  \item Report AIC / BIC for model comparison
  \item Keep all theoretically justified variables
\end{checklist}
\end{multicols}
\end{tcolorbox}

\newpage

% ══════════════════════════════════════════════════════════════════════════════
\section{Qualitative Regressors \& Interaction Effects}
% ══════════════════════════════════════════════════════════════════════════════

\begin{multicols}{2}

\begin{tcolorbox}[when]
\begin{itemize}
  \item Categorical variables (gender, region, industry)
  \item Testing if slopes differ across groups
  \item Policy variables (treatment/control indicators)
  \item Chow test for structural stability
\end{itemize}
\end{tcolorbox}

\begin{tcolorbox}[key]
\textbf{Dummy variable trap:} include $G-1$ dummies for $G$ categories.\\[4pt]
\textbf{Interaction:} $y = \beta_0 + \beta_1 d + \beta_2 x + \beta_3 d \cdot x + u$\\
$\beta_3$ = difference in slope between groups.\\[4pt]
\textbf{Chow test:} $F = \frac{(RSS_R - RSS_U)/k}{RSS_U/(n - 2k)}$
\end{tcolorbox}

\columnbreak

\begin{tcolorbox}[watchout]
\begin{itemize}
  \item Always omit one category (base group)
  \item Interaction terms amplify multicollinearity
  \item Centre continuous vars before interacting
  \item Chow test assumes homoscedasticity
  \item Saturated models with many dummies: overfit risk
\end{itemize}
\end{tcolorbox}

\end{multicols}

\begin{tcolorbox}[checklist]
\begin{multicols}{2}
\begin{checklist}
  \item Choose interpretable base/reference category
  \item Include $G-1$ dummies; verify no perfect collinearity
  \item Test significance of dummy coefficient
  \item Test interaction term if slopes may differ
  \item Chow test if structural break suspected
  \item Report marginal effect at meaningful base group
  \item Plot fitted values by group for visual check
  \item Label coefficients clearly in output
\end{checklist}
\end{multicols}
\end{tcolorbox}

% ══════════════════════════════════════════════════════════════════════════════
\section{Endogeneity \& Instrumental Variables}
% ══════════════════════════════════════════════════════════════════════════════

\begin{multicols}{2}

\begin{tcolorbox}[when]
\begin{itemize}
  \item OVB suspected (ability, motivation, quality)
  \item Simultaneity (supply/demand, wages/education)
  \item Measurement error in key regressor
  \item Natural experiment or policy instrument available
  \item Durbin--Wu--Hausman test rejects exogeneity
\end{itemize}
\end{tcolorbox}

\begin{tcolorbox}[key]
\textbf{OVB formula:} $\text{plim}(\tilde\beta_1) = \beta_1 + \beta_2 \delta_1$\\
where $\delta_1 = \text{Cov}(x_2, x_1)/\text{Var}(x_1)$.\\[4pt]
\textbf{2SLS estimator:}
\[
\hat{\boldsymbol\beta}_{2SLS} = (\hat{\mathbf{X}}'\mathbf{X})^{-1}\hat{\mathbf{X}}'\mathbf{y}
\]
\textbf{First stage:} $x = \pi_0 + \pi_1 z + v$, need $\pi_1 \neq 0$.
\end{tcolorbox}

\columnbreak

\begin{tcolorbox}[watchout]
\begin{itemize}
  \item Weak instruments ($F_{\text{1st stage}} < 10$) inflate 2SLS bias
  \item Instrument must be exogenous (untestable with one IV)
  \item 2SLS SEs larger than OLS by construction
  \item Overidentification: test exclusion with Sargan/Hansen
  \item LATE $\neq$ ATE: IV estimates effect for compliers only
  \item 2SLS inconsistent if instrument invalid
\end{itemize}
\end{tcolorbox}

\end{multicols}

\begin{tcolorbox}[tests]
\begin{tabularx}{\linewidth}{lXl}
\toprule
\textbf{Test} & \textbf{Purpose} & \textbf{Decision} \\
\midrule
First-stage $F$ & Instrument strength & $F > 10$ (Stock--Yogo rule) \\
\rowcolor{RowGray}
Durbin--Wu--Hausman & Endogeneity of regressor & $p < 0.05 \Rightarrow$ use IV \\
Sargan/Hansen $J$ & Overidentification (validity) & $p < 0.05 \Rightarrow$ suspect IV \\
\rowcolor{RowGray}
Partial $R^2$ (1st stage) & Instrument relevance & Higher is better \\
\bottomrule
\end{tabularx}
\end{tcolorbox}

\begin{tcolorbox}[checklist]
\begin{multicols}{2}
\begin{checklist}
  \item State economic argument for endogeneity
  \item Identify candidate instruments (theory-based)
  \item Check instrument relevance: 1st-stage $F > 10$
  \item Argue instrument exogeneity (no direct effect on $y$)
  \item Run Durbin--Wu--Hausman test (compare OLS vs 2SLS)
  \item If $>1$ instrument: run Sargan/Hansen $J$-test
  \item Report both OLS and 2SLS side by side
  \item Interpret 2SLS as LATE if instrument is binary
\end{checklist}
\end{multicols}
\end{tcolorbox}

\newpage

% ══════════════════════════════════════════════════════════════════════════════
\section{Panel Data: Fixed Effects, Random Effects \& First Differences}
% ══════════════════════════════════════════════════════════════════════════════

\begin{multicols}{2}

\begin{tcolorbox}[when]
\begin{itemize}
  \item Data: multiple units observed over multiple periods
  \item Control for unobserved unit-level heterogeneity
  \item FE: if unobserved effect correlated with regressors
  \item RE: if unobserved effect uncorrelated (efficiency gain)
  \item FD: if $T = 2$ or non-stationary panel
\end{itemize}
\end{tcolorbox}

\begin{tcolorbox}[key]
\textbf{FE (within estimator):} demean all variables:\\
$\ddot{y}_{it} = \ddot{\mathbf{x}}_{it}\boldsymbol\beta + \ddot{u}_{it}$\\[4pt]
\textbf{RE:} GLS with $\lambda = 1 - \sigma_u/\sqrt{\sigma_u^2 + T\sigma_a^2}$\\[4pt]
\textbf{FD:} $\Delta y_{it} = \Delta\mathbf{x}_{it}\boldsymbol\beta + \Delta u_{it}$
\end{tcolorbox}

\columnbreak

\begin{tcolorbox}[watchout]
\begin{itemize}
  \item FE removes all time-invariant variables (including gender, race)
  \item FE amplifies attenuation bias from meas.\ error
  \item RE inconsistent if $\text{Cov}(a_i, x_{it}) \neq 0$
  \item Cluster SEs by unit if serial correlation in $u_{it}$
  \item FD loses one time period; FE more efficient with $T > 2$
  \item Time dummies needed to control aggregate shocks
\end{itemize}
\end{tcolorbox}

\end{multicols}

\begin{tcolorbox}[tests]
\begin{tabularx}{\linewidth}{lXl}
\toprule
\textbf{Test} & \textbf{$H_0$} & \textbf{Decision} \\
\midrule
Hausman test & RE consistent ($\alpha_i \perp x_{it}$) & $p < 0.05 \Rightarrow$ use FE \\
\rowcolor{RowGray}
F-test for FE & All $\alpha_i = 0$ (pooled OLS OK) & $p < 0.05 \Rightarrow$ use FE \\
Breusch--Pagan LM & Var($\alpha_i$) $= 0$ (pooled OK) & $p < 0.05 \Rightarrow$ use RE \\
\rowcolor{RowGray}
Serial corr.\ (Wooldridge) & No AR(1) in $u_{it}$ & $p < 0.05 \Rightarrow$ cluster SEs \\
\bottomrule
\end{tabularx}
\end{tcolorbox}

\begin{tcolorbox}[checklist]
\begin{multicols}{2}
\begin{checklist}
  \item Identify panel structure: balanced or unbalanced?
  \item Run pooled OLS as baseline
  \item Run Breusch--Pagan LM test (pool vs RE)
  \item Run FE and RE; save both sets of estimates
  \item Run Hausman test: $p < 0.05 \to$ FE
  \item Add time dummies (test jointly with $F$-test)
  \item Cluster SEs by unit (panel robust)
  \item FE cannot identify time-invariant vars: plan accordingly
\end{checklist}
\end{multicols}
\end{tcolorbox}

% ══════════════════════════════════════════════════════════════════════════════
\section{Stationary Time Series}
% ══════════════════════════════════════════════════════════════════════════════

\begin{multicols}{2}

\begin{tcolorbox}[when]
\begin{itemize}
  \item Confirmed stationary series (ADF rejects unit root)
  \item Forecasting or modelling dynamic relationships
  \item Distributed lag models (PDL, ADL)
  \item Causal analysis between time series
\end{itemize}
\end{tcolorbox}

\begin{tcolorbox}[key]
\textbf{ADL($p$,$q$)} model:\\
$y_t = \alpha + \sum_{j=1}^p \rho_j y_{t-j} + \sum_{s=0}^q \delta_s x_{t-s} + u_t$\\[4pt]
\textbf{Long-run multiplier:} $\frac{\sum \delta_s}{1 - \sum \rho_j}$\\[4pt]
\textbf{HAC (Newey--West) SEs} correct for serial correlation.
\end{tcolorbox}

\columnbreak

\begin{tcolorbox}[watchout]
\begin{itemize}
  \item Trend-stationary vs.\ difference-stationary: different treatment
  \item Seasonality: add seasonal dummies or use seasonal differencing
  \item Spurious regression if levels used on I(1) series
  \item Lag selection: AIC/BIC; don't over-fit
  \item Structural breaks violate stationarity assumption
  \item Granger causality $\neq$ economic causality
\end{itemize}
\end{tcolorbox}

\end{multicols}

\begin{tcolorbox}[checklist]
\begin{multicols}{2}
\begin{checklist}
  \item Plot time series; look for trend, seasonality
  \item Run ADF test (with and without trend)
  \item If I(0): proceed with levels
  \item Include trend/seasonal dummies if needed
  \item Select lag length via AIC or BIC
  \item Use Newey--West (HAC) SEs for inference
  \item Check residuals for remaining autocorrelation (BG)
  \item Granger causality test if dynamic relationship of interest
\end{checklist}
\end{multicols}
\end{tcolorbox}

\newpage

% ══════════════════════════════════════════════════════════════════════════════
\section{Serial Correlation}
% ══════════════════════════════════════════════════════════════════════════════

\begin{multicols}{2}

\begin{tcolorbox}[when]
\begin{itemize}
  \item Time series or panel data residuals
  \item OLS residuals show pattern (ACF plot, DW $\neq 2$)
  \item Dynamically misspecified model
  \item Economic data with inertia (GDP, unemployment)
\end{itemize}
\end{tcolorbox}

\begin{tcolorbox}[key]
\textbf{AR(1):} $u_t = \rho u_{t-1} + e_t,\; |\rho| < 1$\\[4pt]
\textbf{DW statistic:} $DW \approx 2(1-\hat\rho)$\\
$DW \approx 2$: no autocorrelation\\[4pt]
\textbf{Cochrane--Orcutt:} transform by $y_t - \hat\rho y_{t-1}$ and re-estimate.
\end{tcolorbox}

\columnbreak

\begin{tcolorbox}[watchout]
\begin{itemize}
  \item OLS unbiased but inefficient; SEs biased downward
  \item DW test invalid with lagged $y$ as regressor (use BG)
  \item Serial correlation often symptom of omitted dynamics
  \item First fix: add lagged $y$ or lagged $x$ before using HAC
  \item In panels: cluster SEs by unit and period
\end{itemize}
\end{tcolorbox}

\end{multicols}

\begin{tcolorbox}[tests]
\begin{tabularx}{\linewidth}{lXl}
\toprule
\textbf{Test} & \textbf{$H_0$} & \textbf{Decision} \\
\midrule
Durbin--Watson & $\rho = 0$ (no AR(1)) & $DW < d_L$ \reject; avoid if lagged $y$ \\
\rowcolor{RowGray}
Breusch--Godfrey & $H_0$: no AR($p$) & $p < 0.05$ \reject \\
ACF/PACF plot & Visual check & Spikes outside $\pm 1.96/\sqrt{T}$ \\
\rowcolor{RowGray}
Ljung--Box $Q$ & No autocorrelation to lag $h$ & $p < 0.05$ \reject \\
\bottomrule
\end{tabularx}
\end{tcolorbox}

\begin{tcolorbox}[checklist]
\begin{multicols}{2}
\begin{checklist}
  \item Plot ACF and PACF of residuals
  \item Run Breusch--Godfrey test (BG, lag 1 and 4)
  \item If dynamic model: BG preferred over DW
  \item First remedy: add dynamic terms (lags)
  \item If persist: use Newey--West HAC SEs (\texttt{cov\_type='HAC'})
  \item Cochrane--Orcutt for GLS if AR(1) confirmed
  \item After remedy: re-run BG to verify
  \item Never use DW with lagged dependent variable
\end{checklist}
\end{multicols}
\end{tcolorbox}

% ══════════════════════════════════════════════════════════════════════════════
\section{Non-Stationary Time Series: Unit Roots \& Cointegration}
% ══════════════════════════════════════════════════════════════════════════════

\begin{multicols}{2}

\begin{tcolorbox}[when]
\begin{itemize}
  \item ADF fails to reject unit root (series is I(1))
  \item Macro series: GDP, prices, exchange rates
  \item Long-run equilibrium relationship between I(1) series
  \item Error correction model for short-run adjustment
\end{itemize}
\end{tcolorbox}

\begin{tcolorbox}[key]
\textbf{ADF} regression:\\
$\Delta y_t = \alpha + \beta t + \gamma y_{t-1} + \sum_{j=1}^p \delta_j \Delta y_{t-j} + u_t$\\
$H_0$: $\gamma = 0$ (unit root).\\[4pt]
\textbf{ECM:} $\Delta y_t = \alpha + \lambda \hat{e}_{t-1} + \beta \Delta x_t + u_t$\\
$\hat\lambda < 0$: error correction toward equilibrium.
\end{tcolorbox}

\columnbreak

\begin{tcolorbox}[watchout]
\begin{itemize}
  \item Spurious regression: high $R^2$, low DW with I(1) series
  \item First-differencing removes long-run information
  \item Cointegration required before ECM is valid
  \item ADF has low power in small samples (PP test as check)
  \item Include intercept/trend in ADF based on data plot
  \item Johansen for multivariate cointegration
\end{itemize}
\end{tcolorbox}

\end{multicols}

\begin{tcolorbox}[tests]
\begin{tabularx}{\linewidth}{lXl}
\toprule
\textbf{Test} & \textbf{$H_0$} & \textbf{Decision} \\
\midrule
ADF & Unit root ($\gamma = 0$) & $p < 0.05 \Rightarrow$ stationary \\
\rowcolor{RowGray}
Phillips--Perron & Unit root (non-parametric) & $p < 0.05 \Rightarrow$ stationary \\
Engle--Granger & No cointegration (resid I(1)) & $p < 0.05 \Rightarrow$ cointegrated \\
\rowcolor{RowGray}
Johansen trace & $\leq r$ cointegrating vecs & Reject if stat $>$ critical value \\
\bottomrule
\end{tabularx}
\end{tcolorbox}

\begin{tcolorbox}[checklist]
\begin{multicols}{2}
\begin{checklist}
  \item Plot each series; visual check for trend/drift
  \item Run ADF (with constant and with trend)
  \item Also run PP test as robustness check
  \item If I(1): check cointegration (Engle--Granger or Johansen)
  \item If cointegrated: estimate ECM; verify $\hat\lambda < 0$
  \item If not cointegrated: work in first differences
  \item Report DW in levels regression (near 0 = spurious)
  \item Never regress I(1) on I(1) without checking cointegration
\end{checklist}
\end{multicols}
\end{tcolorbox}

\newpage

% ══════════════════════════════════════════════════════════════════════════════
\section{Difference-in-Differences \& Pooled Cross-Sections}
% ══════════════════════════════════════════════════════════════════════════════

\begin{multicols}{2}

\begin{tcolorbox}[when]
\begin{itemize}
  \item Natural experiment: policy affects one group, not another
  \item Pre- and post-treatment observations available
  \item Treatment assignment not random but parallel trends plausible
  \item Pooled cross-sections: repeated surveys (not panel)
\end{itemize}
\end{tcolorbox}

\begin{tcolorbox}[key]
\[
\hat\delta_{\text{DiD}} = (\bar y_{T,\text{after}} - \bar y_{T,\text{before}}) - (\bar y_{C,\text{after}} - \bar y_{C,\text{before}})
\]
Regression form:\\
$y = \beta_0 + \beta_1 D_{\text{treat}} + \beta_2 D_{\text{post}} + \delta D_{\text{treat}} \cdot D_{\text{post}} + u$\\
$\delta$ = ATT under parallel trends.
\end{tcolorbox}

\columnbreak

\begin{tcolorbox}[watchout]
\begin{itemize}
  \item Parallel trends is \textbf{untestable} post-treatment (test pre-trends)
  \item Spillover effects violate SUTVA
  \item Staggered adoption: use Callaway--Sant'Anna, not simple DiD
  \item Always cluster SEs at group (treated/control) level
  \item Anticipation effects can bias DiD
  \item Event study plots: check pre-trend coefficients near zero
\end{itemize}
\end{tcolorbox}

\end{multicols}

\begin{tcolorbox}[checklist]
\begin{multicols}{2}
\begin{checklist}
  \item Plot parallel trends in pre-period (visual + event study)
  \item Define treatment/control groups clearly
  \item Verify no anticipation effects before treatment
  \item Run full DiD regression with interaction term
  \item Cluster SEs by group or state
  \item Add controls to improve precision (not for identification)
  \item Check heterogeneous treatment effects
  \item If staggered: use modern DiD estimator
\end{checklist}
\end{multicols}
\end{tcolorbox}

% ══════════════════════════════════════════════════════════════════════════════
\section{Treatment Effects \& Causal Inference}
% ══════════════════════════════════════════════════════════════════════════════

\begin{multicols}{2}

\begin{tcolorbox}[when]
\begin{itemize}
  \item Estimating causal effect of program/policy/treatment
  \item RCT: randomised control trial available
  \item DiD: quasi-experiment with parallel trends
  \item RDD: assignment determined by threshold rule
  \item PSM: observational data, selection on observables
\end{itemize}
\end{tcolorbox}

\begin{tcolorbox}[key]
\textbf{Potential outcomes:} $Y_i = D_i Y_i(1) + (1-D_i) Y_i(0)$\\[3pt]
$\text{ATE} = E[Y(1) - Y(0)]$\\
$\text{ATT} = E[Y(1) - Y(0) \mid D=1]$\\
$\text{LATE} = \frac{\text{Reduced form}}{\text{First stage}}$ (IV for compliers)\\[3pt]
\textbf{Naive bias:} $\hat\tau_{\text{naive}} = \text{ATT} + \underbrace{E[Y(0)\mid D\!=\!1] - E[Y(0)\mid D\!=\!0]}_{\text{selection bias}}$
\end{tcolorbox}

\columnbreak

\begin{tcolorbox}[watchout]
\begin{itemize}
  \item RCT: OLS on indicator gives ATE unbiasedly
  \item PSM: only controls for \textit{observed} confounders
  \item RDD: local estimate near threshold; check manipulation
  \item LATE $\neq$ ATE unless instrument affects all units (strong)
  \item Always report which estimand (ATE/ATT/LATE) you target
  \item Large bandwidth in RDD: more data, more bias
  \item Common support must hold for PSM (check overlap)
\end{itemize}
\end{tcolorbox}

\end{multicols}

\begin{tcolorbox}[tests]
\begin{tabularx}{\linewidth}{lXl}
\toprule
\textbf{Test / Check} & \textbf{Purpose} & \textbf{Decision} \\
\midrule
Balance table & Covariate balance post-randomisation & All $p > 0.05$ \\
\rowcolor{RowGray}
McCrary density test & No manipulation at RDD cutoff & $p > 0.05$ OK \\
Propensity overlap plot & Common support for PSM & Check visually \\
\rowcolor{RowGray}
Placebo cutoff (RDD) & No effect away from threshold & $p > 0.05$ OK \\
Pre-trend test (DiD) & Parallel pre-trends & Coeff near 0 \\
\bottomrule
\end{tabularx}
\end{tcolorbox}

\begin{tcolorbox}[checklist]
\begin{multicols}{2}
\begin{checklist}
  \item State causal question and target estimand (ATE/ATT/LATE)
  \item Choose identification strategy based on data structure
  \item RCT: verify randomisation (balance table)
  \item DiD: test parallel pre-trends (event study)
  \item RDD: McCrary test; vary bandwidth; placebo cutoff
  \item PSM: estimate propensity score; check overlap
  \item IV/LATE: 1st-stage $F > 10$; instrument exogeneity
  \item Report estimand, sample, and assumptions explicitly
\end{checklist}
\end{multicols}
\end{tcolorbox}

\newpage

% ══════════════════════════════════════════════════════════════════════════════
\section{Limited Dependent Variables: Logit, Probit \& Tobit}
% ══════════════════════════════════════════════════════════════════════════════

\begin{multicols}{2}

\begin{tcolorbox}[when]
\begin{itemize}
  \item \textbf{Logit/Probit:} binary outcome $y \in \{0,1\}$
  \item \textbf{Tobit:} censored/truncated outcome ($y \geq 0$, many zeros)
  \item \textbf{Ordered Probit:} ordinal outcome (Likert scales)
  \item \textbf{Multinomial Logit:} unordered multiple categories
\end{itemize}
\end{tcolorbox}

\begin{tcolorbox}[key]
\textbf{Logit:} $P(y=1|\mathbf{x}) = \Lambda(\mathbf{x}\boldsymbol\beta) = \frac{e^{\mathbf{x}\boldsymbol\beta}}{1+e^{\mathbf{x}\boldsymbol\beta}}$\\[4pt]
\textbf{Marginal effect (AME):} $\frac{\partial P}{\partial x_j} = \hat\beta_j \cdot \bar{\lambda}$\\
where $\bar\lambda = $ avg.\ density evaluated at $\hat{\mathbf{x}}\hat{\boldsymbol{\beta}}$.\\[4pt]
\textbf{Probit:} $\Phi$ replaces $\Lambda$; coefficients $\approx \beta_{\text{logit}}/1.6$.
\end{tcolorbox}

\columnbreak

\begin{tcolorbox}[watchout]
\begin{itemize}
  \item LPM (OLS on binary $y$): predicted probs outside [0,1]
  \item Logit/Probit coefficients: direction only, not magnitude
  \item Always compute and report marginal effects (AME/MEM)
  \item Pseudo-$R^2$ (McFadden) not comparable to OLS $R^2$
  \item Tobit assumes normality; check with residual plots
  \item Perfect separation (predictor fully determines $y$): MLE fails
\end{itemize}
\end{tcolorbox}

\end{multicols}

\begin{tcolorbox}[tests]
\begin{tabularx}{\linewidth}{lXl}
\toprule
\textbf{Test} & \textbf{Purpose} & \textbf{Decision} \\
\midrule
LR test & Joint significance of regressors & $p < 0.05$ \reject \\
\rowcolor{RowGray}
Hosmer--Lemeshow & Goodness of fit (calibration) & $p > 0.05$ good fit \\
Wald test & Individual coefficient $= 0$ & $p < 0.05$ significant \\
\rowcolor{RowGray}
AUC / ROC & Discrimination ability & $> 0.7$ acceptable \\
Vuong test & Logit vs.\ Probit model choice & Significant $\to$ prefer one \\
\bottomrule
\end{tabularx}
\end{tcolorbox}

\begin{tcolorbox}[checklist]
\begin{multicols}{2}
\begin{checklist}
  \item Confirm $y$ is binary/limited before logit/probit
  \item Estimate logit and probit; compare signs/significance
  \item Always compute Average Marginal Effects (AME)
  \item Report pseudo-$R^2$ and classification accuracy
  \item Check for perfect or quasi-separation
  \item Hosmer--Lemeshow test for calibration
  \item Plot ROC curve and report AUC
  \item Do not report raw logit/probit coefficients as effects
\end{checklist}
\end{multicols}
\end{tcolorbox}

% ══════════════════════════════════════════════════════════════════════════════
\section{Advanced Time Series: ARIMA, VAR \& Granger Causality}
% ══════════════════════════════════════════════════════════════════════════════

\begin{multicols}{2}

\begin{tcolorbox}[when]
\begin{itemize}
  \item Univariate forecasting with autocorrelated errors: ARIMA
  \item System of mutually dependent time series: VAR
  \item Testing predictive causality between series: Granger
  \item Impulse response analysis for policy simulations
\end{itemize}
\end{tcolorbox}

\begin{tcolorbox}[key]
\textbf{ARMA($p$,$q$):} $y_t = \sum_{j=1}^p \phi_j y_{t-j} + \varepsilon_t + \sum_{i=1}^q \theta_i \varepsilon_{t-i}$\\[4pt]
\textbf{VAR($p$):} $\mathbf{y}_t = \mathbf{A}_1\mathbf{y}_{t-1} + \cdots + \mathbf{A}_p\mathbf{y}_{t-p} + \mathbf{u}_t$\\[4pt]
\textbf{Granger:} $x$ Granger-causes $y$ if lags of $x$ jointly significant in VAR.
\end{tcolorbox}

\columnbreak

\begin{tcolorbox}[watchout]
\begin{itemize}
  \item ARIMA requires stationarity first (difference if I(1))
  \item Model identification: ACF/PACF, then AIC/BIC
  \item VAR: all series must be stationary (or use VECM if cointegrated)
  \item Granger causality is predictive, not structural causality
  \item Impulse responses depend on Cholesky ordering (economic theory)
  \item Over-parameterised VAR: reduce lag length; LASSO-VAR for large systems
\end{itemize}
\end{tcolorbox}

\end{multicols}

\begin{tcolorbox}[checklist]
\begin{multicols}{2}
\begin{checklist}
  \item Confirm stationarity (ADF) before ARMA/VAR
  \item Inspect ACF/PACF for $p$, $q$ guidance
  \item Estimate candidate ARIMA models; compare AIC/BIC
  \item Check residuals: Ljung--Box test, ACF of residuals
  \item VAR: select lag by AIC/BIC/HQ criteria
  \item Test residuals for autocorrelation (Portmanteau)
  \item Granger causality: Wald test on lag block
  \item Report impulse responses with confidence bands
\end{checklist}
\end{multicols}
\end{tcolorbox}

\newpage

% ══════════════════════════════════════════════════════════════════════════════
\section{Simultaneous Equations Models}
% ══════════════════════════════════════════════════════════════════════════════

\begin{multicols}{2}

\begin{tcolorbox}[when]
\begin{itemize}
  \item Endogenous variables determined jointly (supply \& demand)
  \item Each equation has a causal interpretation
  \item IV/2SLS is insufficient (multiple endogenous vars)
  \item Macro structural models; wage--employment systems
\end{itemize}
\end{tcolorbox}

\begin{tcolorbox}[key]
\textbf{Order condition (necessary):}\\
\# excluded exogenous $\geq$ \# endogenous regressors.\\[4pt]
\textbf{Rank condition (sufficient):} rank of restricted coefficient matrix $= G-1$.\\[4pt]
\textbf{2SLS} equation by equation; \textbf{3SLS} for system efficiency.
\end{tcolorbox}

\columnbreak

\begin{tcolorbox}[watchout]
\begin{itemize}
  \item Under-identified equation: cannot be estimated
  \item OLS on structural equation is inconsistent
  \item Just-identified: 2SLS $=$ IV; over-identified: test with $J$-test
  \item 3SLS more efficient but inconsistency spreads across equations
  \item Normalisation rule: assign one variable as LHS per equation
\end{itemize}
\end{tcolorbox}

\end{multicols}

\begin{tcolorbox}[checklist]
\begin{multicols}{2}
\begin{checklist}
  \item Write structural form explicitly (each equation)
  \item Check order condition for each equation
  \item Verify rank condition before estimation
  \item List all exogenous/excluded variables
  \item Estimate each equation by 2SLS
  \item Sargan test for over-identified equations
  \item Compare 2SLS vs.\ OLS to gauge endogeneity bias
  \item Consider 3SLS if system efficiency matters
\end{checklist}
\end{multicols}
\end{tcolorbox}

% ══════════════════════════════════════════════════════════════════════════════
\section{Diagnostics at a Glance --- Master Reference Table}
% ══════════════════════════════════════════════════════════════════════════════

\small
\begin{tabularx}{\linewidth}{llXll}
\toprule
\rowcolor{NavyBlue}
\color{white}\textbf{Problem} &
\color{white}\textbf{Test} &
\color{white}\textbf{$H_0$} &
\color{white}\textbf{Stat} &
\color{white}\textbf{Remedy} \\
\midrule
Heteroscedasticity & Breusch--Pagan & Homoscedastic & $\chi^2$ & HC1 robust SEs \\
\rowcolor{RowGray}
Heteroscedasticity & White & Homoscedastic & $\chi^2$ & HC1 robust SEs \\
Serial correlation & Breusch--Godfrey & No AR($p$) & $F$ / $\chi^2$ & HAC / C--O \\
\rowcolor{RowGray}
Serial correlation & Durbin--Watson & No AR(1) & $DW$ & Add lags / HAC \\
Multicollinearity & VIF & --- & $>10$ & Drop / PCA / Ridge \\
\rowcolor{RowGray}
Multicollinearity & Condition $\kappa$ & --- & $>30$ & Same as above \\
Misspecification & RESET & Correct form & $F$ & Add nonlinear terms \\
\rowcolor{RowGray}
Unit root & ADF & Unit root & $t$ & Difference / ECM \\
No cointegration & Engle--Granger & Not cointegrated & $t$ & First differences \\
\rowcolor{RowGray}
Endogeneity & Wu--Hausman & Exogenous & $\chi^2$ / $F$ & 2SLS \\
Weak instrument & 1st-stage $F$ & Strong instrument & $F < 10$ & Find better IV \\
\rowcolor{RowGray}
Overidentification & Sargan/Hansen & Valid instruments & $\chi^2$ & Check instruments \\
Panel: RE vs FE & Hausman & RE consistent & $\chi^2$ & FE if $p < 0.05$ \\
\rowcolor{RowGray}
Structural break & Chow & No break & $F$ & Split sample / dummies \\
Normality & Jarque--Bera & Normal residuals & $\chi^2_2$ & Check outliers \\
\bottomrule
\end{tabularx}

\normalsize
\newpage

% ══════════════════════════════════════════════════════════════════════════════
\section{Model Assumptions Summary}
% ══════════════════════════════════════════════════════════════════════════════

\begin{center}
\small
\begin{tabularx}{\linewidth}{lXXl}
\toprule
\rowcolor{NavyBlue}
\color{white}\textbf{Assumption} &
\color{white}\textbf{OLS (cross-section)} &
\color{white}\textbf{Time Series OLS} &
\color{white}\textbf{Violated by} \\
\midrule
MLR.1 & Linear in parameters & TS.1 Same & Misspecification \\
\rowcolor{RowGray}
MLR.2 & Random sample & TS.2 Weak dependence & Non-random, $I(1)$ \\
MLR.3 & No perfect collinearity & TS.3 Same & Dummy trap \\
\rowcolor{RowGray}
MLR.4 & Zero conditional mean & TS.4 Strict exogeneity & Endogeneity, OVB \\
MLR.5 & Homoscedasticity & TS.5 Same & Heteroscedasticity \\
\rowcolor{RowGray}
MLR.6 & Normality of $u$ & TS.6 Same & Outliers, small $n$ \\
(Panel) FE.1 & Strict exogeneity given $a_i$ & --- & Lagged dep.\ var.\ in FE \\
\rowcolor{RowGray}
(IV) IV.1 & Instrument relevance: $\pi_1 \neq 0$ & --- & Weak instruments \\
(IV) IV.2 & Instrument exogeneity: $\text{Cov}(z,u)=0$ & --- & Invalid instrument \\
\bottomrule
\end{tabularx}
\end{center}

% ══════════════════════════════════════════════════════════════════════════════
\section{The Researcher's Pre-Submission Checklist}
% ══════════════════════════════════════════════════════════════════════════════

\begin{tcolorbox}[colback=LightBlue, colframe=NavyBlue, arc=3pt, boxrule=0.8pt,
  title={\bfseries\color{white} Before You Submit: Run Through All of These}, fonttitle=\small]
\begin{multicols}{2}
\textbf{Data \& Design}
\begin{checklist}
  \item Document data source and time period
  \item Check for missing values and their pattern
  \item Report sample size and key descriptive stats
  \item State identification strategy clearly
  \item Confirm exogeneity of key regressors
\end{checklist}

\textbf{OLS Baseline}
\begin{checklist}
  \item Run baseline OLS with full controls
  \item Run RESET test (functional form)
  \item Report robust (HC1) SEs alongside OLS SEs
  \item Check $R^2$, adjusted $R^2$, AIC/BIC
  \item Plot residuals vs.\ fitted values
\end{checklist}

\columnbreak

\textbf{Diagnostic Tests}
\begin{checklist}
  \item Breusch--Pagan and White (heteroscedasticity)
  \item Breusch--Godfrey if time series / panel
  \item VIF for each regressor
  \item ADF if time series data
  \item Hausman if panel data
\end{checklist}

\textbf{Robustness \& Reporting}
\begin{checklist}
  \item Run model on subsamples / alternative specs
  \item Report placebo / falsification tests
  \item Show main table with multiple specifications
  \item Discuss direction and magnitude of violations
  \item State which assumptions cannot be tested
\end{checklist}
\end{multicols}
\end{tcolorbox}

\vspace{1em}

% ══════════════════════════════════════════════════════════════════════════════
\section{Quick Python / statsmodels Reference}
% ══════════════════════════════════════════════════════════════════════════════

\small
\begin{tabularx}{\linewidth}{lX}
\toprule
\rowcolor{NavyBlue}
\color{white}\textbf{Task} & \color{white}\textbf{Python snippet} \\
\midrule
OLS & \texttt{smf.ols('y \textasciitilde{} x1 + x2', data=df).fit()} \\
\rowcolor{RowGray}
Robust SEs & \texttt{res.get\_robustcov\_results(cov\_type='HC1')} \\
HAC SEs & \texttt{res.get\_robustcov\_results(cov\_type='HAC', maxlags=4)} \\
\rowcolor{RowGray}
Breusch--Pagan & \texttt{het\_breuschpagan(res.resid, res.model.exog)} \\
White test & \texttt{het\_white(res.resid, res.model.exog)} \\
\rowcolor{RowGray}
Breusch--Godfrey & \texttt{acorr\_breusch\_godfrey(res, nlags=4)} \\
ADF test & \texttt{adfuller(series, autolag='AIC')} \\
\rowcolor{RowGray}
VIF & \texttt{1 / (1 - smf.ols(f'\{x\} \textasciitilde{} ...').fit().rsquared)} \\
2SLS & \texttt{IV2SLS.from\_formula('y \textasciitilde{} 1 + [x \textasciitilde{} z]', data=df).fit()} \\
\rowcolor{RowGray}
FE panel & \texttt{PanelOLS.from\_formula('y \textasciitilde{} x + EntityEffects', data=pdata).fit()} \\
Hausman & \texttt{compare(fe\_res, re\_res)} from \texttt{linearmodels} \\
\rowcolor{RowGray}
Logit AME & \texttt{logit\_res.get\_margeff().summary()} \\
ARIMA & \texttt{ARIMA(y, order=(p,d,q)).fit()} \\
\rowcolor{RowGray}
Granger & \texttt{grangercausalitytests(data[['y','x']], maxlag=4)} \\
\bottomrule
\end{tabularx}

\normalsize

\vfill
\begin{center}
\small\color{gray}
University of Neuch\^{a}tel --- Python for Econometrics Course\\
This document is a teaching aid and does not replace primary textbooks.\\
Wooldridge (2019) \textit{Introductory Econometrics} $\cdot$
Angrist \& Pischke (2009) \textit{Mostly Harmless Econometrics}
\end{center}

\end{document}